%%%%%%%%%%%%%%%%%%%%%%%%%%%%%%%%%%%%%%%%%%%%%%%%%%%%%%%%%%%%
% 14.6 GEOMETRIC FORMULAS
% The Composition of Advanced Mathematics
%%%%%%%%%%%%%%%%%%%%%%%%%%%%%%%%%%%%%%%%%%%%%%%%%%%%%%%%%%%%

\section{Geometric Formulas}
\label{sec:geometric-formulas}

\prerequisite{
Basic algebra, coordinate geometry, Euclidean geometry,
trigonometry, and mathematical notation.
}

\learningobjectives{
By the end of this reference section, the reader should be able to:
\begin{itemize}
    \item recognize standard notation for geometric quantities;
    \item compute distances and midpoints in coordinate geometry;
    \item compute slopes and equations of lines;
    \item use formulas for angles and polygons;
    \item compute perimeter and area of common plane figures;
    \item compute circumference and area of circles;
    \item use arc-length and sector-area formulas;
    \item recognize formulas for regular polygons;
    \item compute surface areas and volumes of common solids;
    \item recognize formulas for prisms, cylinders, pyramids, cones,
          spheres, and frustums;
    \item use the Pythagorean theorem;
    \item recognize special right-triangle ratios;
    \item use distance formulas in two and three dimensions;
    \item compute areas using coordinates;
    \item recognize vector formulas for geometric area and volume;
    \item use equations of circles, spheres, planes, and common conics;
    \item recognize parametric and polar geometric formulas;
    \item distinguish degrees and radians;
    \item use geometric similarity and scaling laws;
    \item recognize formulas for centroids and centers;
    \item understand how dimensions scale area and volume;
    \item avoid common geometric-formula errors.
\end{itemize}
}

%%%%%%%%%%%%%%%%%%%%%%%%%%%%%%%%%%%%%%%%%%%%%%%%%%%%%%%%%%%%
\subsection{Common Geometric Notation}
\label{subsec:common-geometric-notation}
%%%%%%%%%%%%%%%%%%%%%%%%%%%%%%%%%%%%%%%%%%%%%%%%%%%%%%%%%%%%

Frequently used geometric symbols include

\[
\boxed{
A
=
\text{area},
}
\]

\[
\boxed{
P
=
\text{perimeter},
}
\]

\[
\boxed{
C
=
\text{circumference},
}
\]

\[
\boxed{
V
=
\text{volume},
}
\]

\[
\boxed{
S
=
\text{surface area},
}
\]

\[
\boxed{
r
=
\text{radius},
}
\]

\[
\boxed{
d
=
\text{diameter or distance},
}
\]

\[
\boxed{
h
=
\text{height},
}
\]

and

\[
\boxed{
\theta
=
\text{angle}.
}
\]

Notation varies by context, so quantities should be defined when ambiguity is
possible.

%%%%%%%%%%%%%%%%%%%%%%%%%%%%%%%%%%%%%%%%%%%%%%%%%%%%%%%%%%%%
\subsection{Angle Units}
\label{subsec:angle-units-reference}
%%%%%%%%%%%%%%%%%%%%%%%%%%%%%%%%%%%%%%%%%%%%%%%%%%%%%%%%%%%%

Angles are commonly measured in degrees or radians.

A full revolution is

\[
\boxed{
360^\circ
=
2\pi
\text{ radians}.
}
\]

Therefore,

\[
\boxed{
180^\circ
=
\pi
\text{ radians}.
}
\]

%%%%%%%%%%%%%%%%%%%%%%%%%%%%%%%%%%%%%%%%%%%%%%%%%%%%%%%%%%%%
\subsection{Degree-to-Radian Conversion}
\label{subsec:degree-radian-conversion}
%%%%%%%%%%%%%%%%%%%%%%%%%%%%%%%%%%%%%%%%%%%%%%%%%%%%%%%%%%%%

To convert degrees to radians,

\[
\boxed{
\theta_{\mathrm{rad}}
=
\theta_{\mathrm{deg}}
\frac{\pi}{180}.
}
\]

For example,

\[
60^\circ
=
60
\frac{\pi}{180}
=
\boxed{
\frac{\pi}{3}.
}
\]

%%%%%%%%%%%%%%%%%%%%%%%%%%%%%%%%%%%%%%%%%%%%%%%%%%%%%%%%%%%%
\subsection{Radian-to-Degree Conversion}
\label{subsec:radian-degree-conversion}
%%%%%%%%%%%%%%%%%%%%%%%%%%%%%%%%%%%%%%%%%%%%%%%%%%%%%%%%%%%%

To convert radians to degrees,

\[
\boxed{
\theta_{\mathrm{deg}}
=
\theta_{\mathrm{rad}}
\frac{180}{\pi}.
}
\]

For example,

\[
\frac{3\pi}{4}
=
\frac{3\pi}{4}
\frac{180}{\pi}
=
\boxed{
135^\circ.
}
\]

%%%%%%%%%%%%%%%%%%%%%%%%%%%%%%%%%%%%%%%%%%%%%%%%%%%%%%%%%%%%
\subsection{Complementary and Supplementary Angles}
\label{subsec:complementary-supplementary}
%%%%%%%%%%%%%%%%%%%%%%%%%%%%%%%%%%%%%%%%%%%%%%%%%%%%%%%%%%%%

Complementary angles satisfy

\[
\boxed{
\alpha+\beta
=
90^\circ
=
\frac{\pi}{2}.
}
\]

Supplementary angles satisfy

\[
\boxed{
\alpha+\beta
=
180^\circ
=
\pi.
}
\]

%%%%%%%%%%%%%%%%%%%%%%%%%%%%%%%%%%%%%%%%%%%%%%%%%%%%%%%%%%%%
\subsection{Angles Around a Point}
\label{subsec:angles-around-point}
%%%%%%%%%%%%%%%%%%%%%%%%%%%%%%%%%%%%%%%%%%%%%%%%%%%%%%%%%%%%

Angles around a point sum to

\[
\boxed{
360^\circ
=
2\pi.
}
\]

Thus if angles \(\theta_1,\ldots,\theta_n\) meet around a point,

\[
\boxed{
\sum_{k=1}^{n}\theta_k
=
360^\circ.
}
\]

%%%%%%%%%%%%%%%%%%%%%%%%%%%%%%%%%%%%%%%%%%%%%%%%%%%%%%%%%%%%
\subsection{Triangle Angle Sum}
\label{subsec:triangle-angle-sum}
%%%%%%%%%%%%%%%%%%%%%%%%%%%%%%%%%%%%%%%%%%%%%%%%%%%%%%%%%%%%

For a Euclidean triangle,

\[
\boxed{
A+B+C
=
180^\circ
=
\pi.
}
\]

This formula depends on Euclidean geometry.

%%%%%%%%%%%%%%%%%%%%%%%%%%%%%%%%%%%%%%%%%%%%%%%%%%%%%%%%%%%%
\subsection{Exterior Angle of a Triangle}
\label{subsec:triangle-exterior-angle}
%%%%%%%%%%%%%%%%%%%%%%%%%%%%%%%%%%%%%%%%%%%%%%%%%%%%%%%%%%%%

An exterior angle of a triangle equals the sum of the two remote interior
angles.

If exterior angle is \(E\), then

\[
\boxed{
E=A+B.
}
\]

%%%%%%%%%%%%%%%%%%%%%%%%%%%%%%%%%%%%%%%%%%%%%%%%%%%%%%%%%%%%
\subsection{Polygon Interior Angle Sum}
\label{subsec:polygon-interior-angle-sum}
%%%%%%%%%%%%%%%%%%%%%%%%%%%%%%%%%%%%%%%%%%%%%%%%%%%%%%%%%%%%

For a simple \(n\)-gon,

\[
\boxed{
\text{sum of interior angles}
=
(n-2)180^\circ.
}
\]

In radians,

\[
\boxed{
(n-2)\pi.
}
\]

%%%%%%%%%%%%%%%%%%%%%%%%%%%%%%%%%%%%%%%%%%%%%%%%%%%%%%%%%%%%
\subsection{Regular Polygon Interior Angle}
\label{subsec:regular-polygon-interior-angle}
%%%%%%%%%%%%%%%%%%%%%%%%%%%%%%%%%%%%%%%%%%%%%%%%%%%%%%%%%%%%

For a regular \(n\)-gon, each interior angle is

\[
\boxed{
\frac{(n-2)180^\circ}{n}.
}
\]

In radians,

\[
\boxed{
\frac{(n-2)\pi}{n}.
}
\]

%%%%%%%%%%%%%%%%%%%%%%%%%%%%%%%%%%%%%%%%%%%%%%%%%%%%%%%%%%%%
\subsection{Regular Polygon Exterior Angle}
\label{subsec:regular-polygon-exterior-angle}
%%%%%%%%%%%%%%%%%%%%%%%%%%%%%%%%%%%%%%%%%%%%%%%%%%%%%%%%%%%%

Each exterior angle of a regular \(n\)-gon is

\[
\boxed{
\frac{360^\circ}{n}.
}
\]

In radians,

\[
\boxed{
\frac{2\pi}{n}.
}
\]

%%%%%%%%%%%%%%%%%%%%%%%%%%%%%%%%%%%%%%%%%%%%%%%%%%%%%%%%%%%%
\subsection{Pythagorean Theorem}
\label{subsec:pythagorean-theorem-reference}
%%%%%%%%%%%%%%%%%%%%%%%%%%%%%%%%%%%%%%%%%%%%%%%%%%%%%%%%%%%%

For a right triangle with legs \(a,b\) and hypotenuse \(c\),

\[
\boxed{
a^2+b^2=c^2.
}
\]

Therefore,

\[
\boxed{
c
=
\sqrt{a^2+b^2}.
}
\]

%%%%%%%%%%%%%%%%%%%%%%%%%%%%%%%%%%%%%%%%%%%%%%%%%%%%%%%%%%%%
\subsection{Worked Example: Pythagorean Theorem}
\label{subsec:worked-pythagorean-reference}
%%%%%%%%%%%%%%%%%%%%%%%%%%%%%%%%%%%%%%%%%%%%%%%%%%%%%%%%%%%%

\begin{workedexamplebox}{Find the Hypotenuse}

Suppose

\[
a=5,
\qquad
b=12.
\]

Then

\[
c^2
=
5^2+12^2
=
25+144
=
169.
\]

Thus

\begin{answerbox}
\[
\boxed{
c=13.
}
\]
\end{answerbox}

\end{workedexamplebox}

%%%%%%%%%%%%%%%%%%%%%%%%%%%%%%%%%%%%%%%%%%%%%%%%%%%%%%%%%%%%
\subsection{Special Right Triangle: \(45^\circ-45^\circ-90^\circ\)}
\label{subsec:45-45-90-reference}
%%%%%%%%%%%%%%%%%%%%%%%%%%%%%%%%%%%%%%%%%%%%%%%%%%%%%%%%%%%%

The side ratio is

\[
\boxed{
1:1:\sqrt2.
}
\]

If each leg has length \(a\), then the hypotenuse is

\[
\boxed{
a\sqrt2.
}
\]

%%%%%%%%%%%%%%%%%%%%%%%%%%%%%%%%%%%%%%%%%%%%%%%%%%%%%%%%%%%%
\subsection{Special Right Triangle: \(30^\circ-60^\circ-90^\circ\)}
\label{subsec:30-60-90-reference}
%%%%%%%%%%%%%%%%%%%%%%%%%%%%%%%%%%%%%%%%%%%%%%%%%%%%%%%%%%%%

The side ratio is

\[
\boxed{
1:\sqrt3:2.
}
\]

If the shortest side is \(a\), then

\[
\boxed{
\text{longer leg}
=
a\sqrt3
}
\]

and

\[
\boxed{
\text{hypotenuse}
=
2a.
}
\]

%%%%%%%%%%%%%%%%%%%%%%%%%%%%%%%%%%%%%%%%%%%%%%%%%%%%%%%%%%%%
\subsection{Triangle Perimeter}
\label{subsec:triangle-perimeter-reference}
%%%%%%%%%%%%%%%%%%%%%%%%%%%%%%%%%%%%%%%%%%%%%%%%%%%%%%%%%%%%

For side lengths \(a,b,c\),

\[
\boxed{
P=a+b+c.
}
\]

%%%%%%%%%%%%%%%%%%%%%%%%%%%%%%%%%%%%%%%%%%%%%%%%%%%%%%%%%%%%
\subsection{Triangle Area: Base and Height}
\label{subsec:triangle-area-base-height}
%%%%%%%%%%%%%%%%%%%%%%%%%%%%%%%%%%%%%%%%%%%%%%%%%%%%%%%%%%%%

For base \(b\) and perpendicular height \(h\),

\[
\boxed{
A
=
\frac12bh.
}
\]

%%%%%%%%%%%%%%%%%%%%%%%%%%%%%%%%%%%%%%%%%%%%%%%%%%%%%%%%%%%%
\subsection{Heron's Formula}
\label{subsec:herons-formula-reference}
%%%%%%%%%%%%%%%%%%%%%%%%%%%%%%%%%%%%%%%%%%%%%%%%%%%%%%%%%%%%

For triangle sides \(a,b,c\), define semiperimeter

\[
\boxed{
s
=
\frac{a+b+c}{2}.
}
\]

Then

\[
\boxed{
A
=
\sqrt{
s(s-a)(s-b)(s-c)
}.
}
\]

%%%%%%%%%%%%%%%%%%%%%%%%%%%%%%%%%%%%%%%%%%%%%%%%%%%%%%%%%%%%
\subsection{Triangle Area From Two Sides and an Included Angle}
\label{subsec:triangle-area-trig}
%%%%%%%%%%%%%%%%%%%%%%%%%%%%%%%%%%%%%%%%%%%%%%%%%%%%%%%%%%%%

If sides \(a,b\) include angle \(C\),

\[
\boxed{
A
=
\frac12ab\sin C.
}
\]

Similarly,

\[
\boxed{
A
=
\frac12bc\sin A
=
\frac12ca\sin B.
}
\]

%%%%%%%%%%%%%%%%%%%%%%%%%%%%%%%%%%%%%%%%%%%%%%%%%%%%%%%%%%%%
\subsection{Equilateral Triangle}
\label{subsec:equilateral-triangle-formulas}
%%%%%%%%%%%%%%%%%%%%%%%%%%%%%%%%%%%%%%%%%%%%%%%%%%%%%%%%%%%%

For side length \(a\), height is

\[
\boxed{
h
=
\frac{\sqrt3}{2}a.
}
\]

Area is

\[
\boxed{
A
=
\frac{\sqrt3}{4}a^2.
}
\]

Perimeter is

\[
\boxed{
P=3a.
}
\]

%%%%%%%%%%%%%%%%%%%%%%%%%%%%%%%%%%%%%%%%%%%%%%%%%%%%%%%%%%%%
\subsection{Rectangle}
\label{subsec:rectangle-formulas}
%%%%%%%%%%%%%%%%%%%%%%%%%%%%%%%%%%%%%%%%%%%%%%%%%%%%%%%%%%%%

For length \(l\) and width \(w\),

\[
\boxed{
A=lw.
}
\]

Perimeter is

\[
\boxed{
P=2l+2w
=
2(l+w).
}
\]

Diagonal length is

\[
\boxed{
d
=
\sqrt{l^2+w^2}.
}
\]

%%%%%%%%%%%%%%%%%%%%%%%%%%%%%%%%%%%%%%%%%%%%%%%%%%%%%%%%%%%%
\subsection{Square}
\label{subsec:square-formulas}
%%%%%%%%%%%%%%%%%%%%%%%%%%%%%%%%%%%%%%%%%%%%%%%%%%%%%%%%%%%%

For side \(s\),

\[
\boxed{
A=s^2,
}
\]

\[
\boxed{
P=4s,
}
\]

and diagonal

\[
\boxed{
d=s\sqrt2.
}
\]

%%%%%%%%%%%%%%%%%%%%%%%%%%%%%%%%%%%%%%%%%%%%%%%%%%%%%%%%%%%%
\subsection{Parallelogram}
\label{subsec:parallelogram-formulas}
%%%%%%%%%%%%%%%%%%%%%%%%%%%%%%%%%%%%%%%%%%%%%%%%%%%%%%%%%%%%

For base \(b\) and perpendicular height \(h\),

\[
\boxed{
A=bh.
}
\]

For adjacent sides \(a,b\) with included angle \(\theta\),

\[
\boxed{
A
=
ab\sin\theta.
}
\]

%%%%%%%%%%%%%%%%%%%%%%%%%%%%%%%%%%%%%%%%%%%%%%%%%%%%%%%%%%%%
\subsection{Rhombus}
\label{subsec:rhombus-formulas}
%%%%%%%%%%%%%%%%%%%%%%%%%%%%%%%%%%%%%%%%%%%%%%%%%%%%%%%%%%%%

For diagonals \(d_1,d_2\),

\[
\boxed{
A
=
\frac12d_1d_2.
}
\]

For side \(a\) and included angle \(\theta\),

\[
\boxed{
A
=
a^2\sin\theta.
}
\]

%%%%%%%%%%%%%%%%%%%%%%%%%%%%%%%%%%%%%%%%%%%%%%%%%%%%%%%%%%%%
\subsection{Trapezoid}
\label{subsec:trapezoid-formulas}
%%%%%%%%%%%%%%%%%%%%%%%%%%%%%%%%%%%%%%%%%%%%%%%%%%%%%%%%%%%%

For parallel bases \(b_1,b_2\) and height \(h\),

\[
\boxed{
A
=
\frac12
(b_1+b_2)h.
}
\]

The median or midsegment length is

\[
\boxed{
m
=
\frac{b_1+b_2}{2}.
}
\]

Thus

\[
\boxed{
A=mh.
}
\]

%%%%%%%%%%%%%%%%%%%%%%%%%%%%%%%%%%%%%%%%%%%%%%%%%%%%%%%%%%%%
\subsection{Kite}
\label{subsec:kite-formulas}
%%%%%%%%%%%%%%%%%%%%%%%%%%%%%%%%%%%%%%%%%%%%%%%%%%%%%%%%%%%%

For perpendicular diagonals \(d_1,d_2\),

\[
\boxed{
A
=
\frac12d_1d_2.
}
\]

%%%%%%%%%%%%%%%%%%%%%%%%%%%%%%%%%%%%%%%%%%%%%%%%%%%%%%%%%%%%
\subsection{Regular Polygon Perimeter}
\label{subsec:regular-polygon-perimeter}
%%%%%%%%%%%%%%%%%%%%%%%%%%%%%%%%%%%%%%%%%%%%%%%%%%%%%%%%%%%%

For a regular \(n\)-gon with side length \(s\),

\[
\boxed{
P=ns.
}
\]

%%%%%%%%%%%%%%%%%%%%%%%%%%%%%%%%%%%%%%%%%%%%%%%%%%%%%%%%%%%%
\subsection{Regular Polygon Area Using Apothem}
\label{subsec:regular-polygon-area-apothem}
%%%%%%%%%%%%%%%%%%%%%%%%%%%%%%%%%%%%%%%%%%%%%%%%%%%%%%%%%%%%

If \(a\) is the apothem and \(P\) is perimeter,

\[
\boxed{
A
=
\frac12aP.
}
\]

Thus for side length \(s\),

\[
\boxed{
A
=
\frac12a(ns).
}
\]

%%%%%%%%%%%%%%%%%%%%%%%%%%%%%%%%%%%%%%%%%%%%%%%%%%%%%%%%%%%%
\subsection{Regular Polygon Area From Side Length}
\label{subsec:regular-polygon-area-side}
%%%%%%%%%%%%%%%%%%%%%%%%%%%%%%%%%%%%%%%%%%%%%%%%%%%%%%%%%%%%

For a regular \(n\)-gon with side length \(s\),

\[
\boxed{
A
=
\frac{
ns^2
}{
4\tan(\pi/n)
}.
}
\]

%%%%%%%%%%%%%%%%%%%%%%%%%%%%%%%%%%%%%%%%%%%%%%%%%%%%%%%%%%%%
\subsection{Circle Diameter and Radius}
\label{subsec:circle-radius-diameter}
%%%%%%%%%%%%%%%%%%%%%%%%%%%%%%%%%%%%%%%%%%%%%%%%%%%%%%%%%%%%

For a circle,

\[
\boxed{
d=2r.
}
\]

Therefore,

\[
\boxed{
r=\frac{d}{2}.
}
\]

%%%%%%%%%%%%%%%%%%%%%%%%%%%%%%%%%%%%%%%%%%%%%%%%%%%%%%%%%%%%
\subsection{Circle Circumference}
\label{subsec:circle-circumference-reference}
%%%%%%%%%%%%%%%%%%%%%%%%%%%%%%%%%%%%%%%%%%%%%%%%%%%%%%%%%%%%

\[
\boxed{
C
=
2\pi r
}
\]

or

\[
\boxed{
C
=
\pi d.
}
\]

%%%%%%%%%%%%%%%%%%%%%%%%%%%%%%%%%%%%%%%%%%%%%%%%%%%%%%%%%%%%
\subsection{Circle Area}
\label{subsec:circle-area-reference}
%%%%%%%%%%%%%%%%%%%%%%%%%%%%%%%%%%%%%%%%%%%%%%%%%%%%%%%%%%%%

\[
\boxed{
A
=
\pi r^2.
}
\]

%%%%%%%%%%%%%%%%%%%%%%%%%%%%%%%%%%%%%%%%%%%%%%%%%%%%%%%%%%%%
\subsection{Arc Length}
\label{subsec:arc-length-circle-reference}
%%%%%%%%%%%%%%%%%%%%%%%%%%%%%%%%%%%%%%%%%%%%%%%%%%%%%%%%%%%%

For central angle \(\theta\) measured in radians,

\[
\boxed{
s
=
r\theta.
}
\]

If \(\theta\) is in degrees,

\[
\boxed{
s
=
\frac{\theta}{360^\circ}
(2\pi r).
}
\]

%%%%%%%%%%%%%%%%%%%%%%%%%%%%%%%%%%%%%%%%%%%%%%%%%%%%%%%%%%%%
\subsection{Sector Area}
\label{subsec:sector-area-reference}
%%%%%%%%%%%%%%%%%%%%%%%%%%%%%%%%%%%%%%%%%%%%%%%%%%%%%%%%%%%%

For \(\theta\) in radians,

\[
\boxed{
A_{\mathrm{sector}}
=
\frac12r^2\theta.
}
\]

For degrees,

\[
\boxed{
A_{\mathrm{sector}}
=
\frac{\theta}{360^\circ}
\pi r^2.
}
\]

%%%%%%%%%%%%%%%%%%%%%%%%%%%%%%%%%%%%%%%%%%%%%%%%%%%%%%%%%%%%
\subsection{Chord Length}
\label{subsec:chord-length-reference}
%%%%%%%%%%%%%%%%%%%%%%%%%%%%%%%%%%%%%%%%%%%%%%%%%%%%%%%%%%%%

For a circle of radius \(r\) and central angle \(\theta\),

\[
\boxed{
c
=
2r
\sin
\left(
\frac{\theta}{2}
\right).
}
\]

Here \(\theta\) may be interpreted consistently in either angular unit when
used inside the corresponding trigonometric function convention.

%%%%%%%%%%%%%%%%%%%%%%%%%%%%%%%%%%%%%%%%%%%%%%%%%%%%%%%%%%%%
\subsection{Segment Area}
\label{subsec:circular-segment-area}
%%%%%%%%%%%%%%%%%%%%%%%%%%%%%%%%%%%%%%%%%%%%%%%%%%%%%%%%%%%%

For central angle \(\theta\) in radians,

\[
\boxed{
A_{\mathrm{segment}}
=
\frac12r^2
(\theta-\sin\theta).
}
\]

This equals sector area minus triangular area.

%%%%%%%%%%%%%%%%%%%%%%%%%%%%%%%%%%%%%%%%%%%%%%%%%%%%%%%%%%%%
\subsection{Annulus}
\label{subsec:annulus-formula}
%%%%%%%%%%%%%%%%%%%%%%%%%%%%%%%%%%%%%%%%%%%%%%%%%%%%%%%%%%%%

For outer radius \(R\) and inner radius \(r\),

\[
\boxed{
A
=
\pi
(R^2-r^2).
}
\]

%%%%%%%%%%%%%%%%%%%%%%%%%%%%%%%%%%%%%%%%%%%%%%%%%%%%%%%%%%%%
\subsection{Ellipse}
\label{subsec:ellipse-area-reference}
%%%%%%%%%%%%%%%%%%%%%%%%%%%%%%%%%%%%%%%%%%%%%%%%%%%%%%%%%%%%

For semimajor axis \(a\) and semiminor axis \(b\),

\[
\boxed{
A
=
\pi ab.
}
\]

The ellipse

\[
\boxed{
\frac{x^2}{a^2}
+
\frac{y^2}{b^2}
=
1
}
\]

has semiaxes \(a\) and \(b\).

%%%%%%%%%%%%%%%%%%%%%%%%%%%%%%%%%%%%%%%%%%%%%%%%%%%%%%%%%%%%
\subsection{Ellipse Eccentricity}
\label{subsec:ellipse-eccentricity-reference}
%%%%%%%%%%%%%%%%%%%%%%%%%%%%%%%%%%%%%%%%%%%%%%%%%%%%%%%%%%%%

For \(a\geq b>0\),

\[
\boxed{
c
=
\sqrt{a^2-b^2}
}
\]

is the focal distance from center.

Eccentricity is

\[
\boxed{
e
=
\frac{c}{a}.
}
\]

Thus

\[
\boxed{
0\leq e<1.
}
\]

%%%%%%%%%%%%%%%%%%%%%%%%%%%%%%%%%%%%%%%%%%%%%%%%%%%%%%%%%%%%
\subsection{Parabola Standard Form}
\label{subsec:parabola-standard-reference}
%%%%%%%%%%%%%%%%%%%%%%%%%%%%%%%%%%%%%%%%%%%%%%%%%%%%%%%%%%%%

A horizontal-axis parabola may be written

\[
\boxed{
(y-k)^2
=
4p(x-h).
}
\]

Its vertex is

\[
\boxed{
(h,k).
}
\]

Its focus is

\[
\boxed{
(h+p,k).
}
\]

Its directrix is

\[
\boxed{
x=h-p.
}
\]

%%%%%%%%%%%%%%%%%%%%%%%%%%%%%%%%%%%%%%%%%%%%%%%%%%%%%%%%%%%%
\subsection{Vertical Parabola Standard Form}
\label{subsec:vertical-parabola-standard}
%%%%%%%%%%%%%%%%%%%%%%%%%%%%%%%%%%%%%%%%%%%%%%%%%%%%%%%%%%%%

A vertical parabola may be written

\[
\boxed{
(x-h)^2
=
4p(y-k).
}
\]

Its focus is

\[
\boxed{
(h,k+p),
}
\]

and directrix is

\[
\boxed{
y=k-p.
}
\]

%%%%%%%%%%%%%%%%%%%%%%%%%%%%%%%%%%%%%%%%%%%%%%%%%%%%%%%%%%%%
\subsection{Hyperbola Standard Form}
\label{subsec:hyperbola-standard-reference}
%%%%%%%%%%%%%%%%%%%%%%%%%%%%%%%%%%%%%%%%%%%%%%%%%%%%%%%%%%%%

A horizontal hyperbola has form

\[
\boxed{
\frac{(x-h)^2}{a^2}
-
\frac{(y-k)^2}{b^2}
=
1.
}
\]

Its asymptotes are

\[
\boxed{
y-k
=
\pm
\frac{b}{a}
(x-h).
}
\]

%%%%%%%%%%%%%%%%%%%%%%%%%%%%%%%%%%%%%%%%%%%%%%%%%%%%%%%%%%%%
\subsection{Hyperbola Focal Relation}
\label{subsec:hyperbola-focal-relation}
%%%%%%%%%%%%%%%%%%%%%%%%%%%%%%%%%%%%%%%%%%%%%%%%%%%%%%%%%%%%

For a hyperbola,

\[
\boxed{
c^2
=
a^2+b^2.
}
\]

Its eccentricity is

\[
\boxed{
e
=
\frac{c}{a}
>
1.
}
\]

%%%%%%%%%%%%%%%%%%%%%%%%%%%%%%%%%%%%%%%%%%%%%%%%%%%%%%%%%%%%
\subsection{Coordinate Distance Formula}
\label{subsec:coordinate-distance-formula}
%%%%%%%%%%%%%%%%%%%%%%%%%%%%%%%%%%%%%%%%%%%%%%%%%%%%%%%%%%%%

For points

\[
P_1=(x_1,y_1),
\qquad
P_2=(x_2,y_2),
\]

distance is

\[
\boxed{
d
=
\sqrt{
(x_2-x_1)^2
+
(y_2-y_1)^2
}.
}
\]

%%%%%%%%%%%%%%%%%%%%%%%%%%%%%%%%%%%%%%%%%%%%%%%%%%%%%%%%%%%%
\subsection{Three-Dimensional Distance}
\label{subsec:three-dimensional-distance}
%%%%%%%%%%%%%%%%%%%%%%%%%%%%%%%%%%%%%%%%%%%%%%%%%%%%%%%%%%%%

For

\[
P_1=(x_1,y_1,z_1)
\]

and

\[
P_2=(x_2,y_2,z_2),
\]

\[
\boxed{
d
=
\sqrt{
(x_2-x_1)^2
+
(y_2-y_1)^2
+
(z_2-z_1)^2
}.
}
\]

%%%%%%%%%%%%%%%%%%%%%%%%%%%%%%%%%%%%%%%%%%%%%%%%%%%%%%%%%%%%
\subsection{Midpoint Formula}
\label{subsec:midpoint-formula-reference}
%%%%%%%%%%%%%%%%%%%%%%%%%%%%%%%%%%%%%%%%%%%%%%%%%%%%%%%%%%%%

For endpoints

\[
(x_1,y_1),
\qquad
(x_2,y_2),
\]

the midpoint is

\[
\boxed{
M
=
\left(
\frac{x_1+x_2}{2},
\frac{y_1+y_2}{2}
\right).
}
\]

In three dimensions,

\[
\boxed{
M
=
\left(
\frac{x_1+x_2}{2},
\frac{y_1+y_2}{2},
\frac{z_1+z_2}{2}
\right).
}
\]

%%%%%%%%%%%%%%%%%%%%%%%%%%%%%%%%%%%%%%%%%%%%%%%%%%%%%%%%%%%%
\subsection{Slope}
\label{subsec:slope-reference}
%%%%%%%%%%%%%%%%%%%%%%%%%%%%%%%%%%%%%%%%%%%%%%%%%%%%%%%%%%%%

For distinct points with \(x_2\neq x_1\),

\[
\boxed{
m
=
\frac{
y_2-y_1
}{
x_2-x_1
}.
}
\]

Slope measures vertical change divided by horizontal change.

%%%%%%%%%%%%%%%%%%%%%%%%%%%%%%%%%%%%%%%%%%%%%%%%%%%%%%%%%%%%
\subsection{Slope-Intercept Form}
\label{subsec:slope-intercept-reference}
%%%%%%%%%%%%%%%%%%%%%%%%%%%%%%%%%%%%%%%%%%%%%%%%%%%%%%%%%%%%

A nonvertical line may be written

\[
\boxed{
y=mx+b.
}
\]

Here

\[
m
=
\text{slope}
\]

and

\[
b
=
\text{\(y\)-intercept}.
\]

%%%%%%%%%%%%%%%%%%%%%%%%%%%%%%%%%%%%%%%%%%%%%%%%%%%%%%%%%%%%
\subsection{Point-Slope Form}
\label{subsec:point-slope-reference}
%%%%%%%%%%%%%%%%%%%%%%%%%%%%%%%%%%%%%%%%%%%%%%%%%%%%%%%%%%%%

A line with slope \(m\) passing through \((x_1,y_1)\) satisfies

\[
\boxed{
y-y_1
=
m(x-x_1).
}
\]

%%%%%%%%%%%%%%%%%%%%%%%%%%%%%%%%%%%%%%%%%%%%%%%%%%%%%%%%%%%%
\subsection{Two-Point Form}
\label{subsec:two-point-line-reference}
%%%%%%%%%%%%%%%%%%%%%%%%%%%%%%%%%%%%%%%%%%%%%%%%%%%%%%%%%%%%

For \(x_1\neq x_2\),

\[
\boxed{
y-y_1
=
\frac{
y_2-y_1
}{
x_2-x_1
}
(x-x_1).
}
\]

%%%%%%%%%%%%%%%%%%%%%%%%%%%%%%%%%%%%%%%%%%%%%%%%%%%%%%%%%%%%
\subsection{Standard Line Form}
\label{subsec:standard-line-reference}
%%%%%%%%%%%%%%%%%%%%%%%%%%%%%%%%%%%%%%%%%%%%%%%%%%%%%%%%%%%%

A line may be written

\[
\boxed{
Ax+By=C.
}
\]

If \(B\neq0\), then slope is

\[
\boxed{
m
=
-\frac{A}{B}.
}
\]

%%%%%%%%%%%%%%%%%%%%%%%%%%%%%%%%%%%%%%%%%%%%%%%%%%%%%%%%%%%%
\subsection{Parallel Lines}
\label{subsec:parallel-lines-reference}
%%%%%%%%%%%%%%%%%%%%%%%%%%%%%%%%%%%%%%%%%%%%%%%%%%%%%%%%%%%%

Two distinct nonvertical lines are parallel when

\[
\boxed{
m_1=m_2.
}
\]

In vector form, their direction vectors are scalar multiples.

%%%%%%%%%%%%%%%%%%%%%%%%%%%%%%%%%%%%%%%%%%%%%%%%%%%%%%%%%%%%
\subsection{Perpendicular Lines}
\label{subsec:perpendicular-lines-reference}
%%%%%%%%%%%%%%%%%%%%%%%%%%%%%%%%%%%%%%%%%%%%%%%%%%%%%%%%%%%%

For nonvertical, nonhorizontal lines,

\[
\boxed{
m_1m_2=-1.
}
\]

Thus their slopes are negative reciprocals.

%%%%%%%%%%%%%%%%%%%%%%%%%%%%%%%%%%%%%%%%%%%%%%%%%%%%%%%%%%%%
\subsection{Distance From a Point to a Line}
\label{subsec:point-line-distance}
%%%%%%%%%%%%%%%%%%%%%%%%%%%%%%%%%%%%%%%%%%%%%%%%%%%%%%%%%%%%

For point \((x_0,y_0)\) and line

\[
Ax+By+C=0,
\]

distance is

\[
\boxed{
d
=
\frac{
|Ax_0+By_0+C|
}{
\sqrt{A^2+B^2}
}.
}
\]

%%%%%%%%%%%%%%%%%%%%%%%%%%%%%%%%%%%%%%%%%%%%%%%%%%%%%%%%%%%%
\subsection{Equation of a Circle}
\label{subsec:circle-equation-reference}
%%%%%%%%%%%%%%%%%%%%%%%%%%%%%%%%%%%%%%%%%%%%%%%%%%%%%%%%%%%%

A circle with center

\[
(h,k)
\]

and radius \(r\) satisfies

\[
\boxed{
(x-h)^2
+
(y-k)^2
=
r^2.
}
\]

At the origin,

\[
\boxed{
x^2+y^2=r^2.
}
\]

%%%%%%%%%%%%%%%%%%%%%%%%%%%%%%%%%%%%%%%%%%%%%%%%%%%%%%%%%%%%
\subsection{Equation of a Sphere}
\label{subsec:sphere-equation-reference}
%%%%%%%%%%%%%%%%%%%%%%%%%%%%%%%%%%%%%%%%%%%%%%%%%%%%%%%%%%%%

A sphere centered at

\[
(h,k,\ell)
\]

with radius \(r\) satisfies

\[
\boxed{
(x-h)^2
+
(y-k)^2
+
(z-\ell)^2
=
r^2.
}
\]

%%%%%%%%%%%%%%%%%%%%%%%%%%%%%%%%%%%%%%%%%%%%%%%%%%%%%%%%%%%%
\subsection{Area of a Triangle From Coordinates}
\label{subsec:coordinate-triangle-area}
%%%%%%%%%%%%%%%%%%%%%%%%%%%%%%%%%%%%%%%%%%%%%%%%%%%%%%%%%%%%

For vertices

\[
(x_1,y_1),
\quad
(x_2,y_2),
\quad
(x_3,y_3),
\]

the area is

\[
\boxed{
A
=
\frac12
\left|
x_1(y_2-y_3)
+
x_2(y_3-y_1)
+
x_3(y_1-y_2)
\right|.
}
\]

%%%%%%%%%%%%%%%%%%%%%%%%%%%%%%%%%%%%%%%%%%%%%%%%%%%%%%%%%%%%
\subsection{Determinant Form of Triangle Area}
\label{subsec:triangle-area-determinant}
%%%%%%%%%%%%%%%%%%%%%%%%%%%%%%%%%%%%%%%%%%%%%%%%%%%%%%%%%%%%

Equivalently,

\[
\boxed{
A
=
\frac12
\left|
\det
\begin{bmatrix}
x_1 & y_1 & 1\\
x_2 & y_2 & 1\\
x_3 & y_3 & 1
\end{bmatrix}
\right|.
}
\]

%%%%%%%%%%%%%%%%%%%%%%%%%%%%%%%%%%%%%%%%%%%%%%%%%%%%%%%%%%%%
\subsection{Shoelace Formula}
\label{subsec:shoelace-formula-reference}
%%%%%%%%%%%%%%%%%%%%%%%%%%%%%%%%%%%%%%%%%%%%%%%%%%%%%%%%%%%%

For a simple polygon with ordered vertices

\[
(x_1,y_1),
\ldots,
(x_n,y_n),
\]

and

\[
(x_{n+1},y_{n+1})
=
(x_1,y_1),
\]

the area is

\[
\boxed{
A
=
\frac12
\left|
\sum_{i=1}^{n}
(x_iy_{i+1}-y_ix_{i+1})
\right|.
}
\]

%%%%%%%%%%%%%%%%%%%%%%%%%%%%%%%%%%%%%%%%%%%%%%%%%%%%%%%%%%%%
\subsection{Vector Length}
\label{subsec:vector-length-geometric-reference}
%%%%%%%%%%%%%%%%%%%%%%%%%%%%%%%%%%%%%%%%%%%%%%%%%%%%%%%%%%%%

For

\[
\mathbf{v}
=
(v_1,\ldots,v_n),
\]

the Euclidean length is

\[
\boxed{
\|\mathbf{v}\|
=
\sqrt{
v_1^2+\cdots+v_n^2
}.
}
\]

%%%%%%%%%%%%%%%%%%%%%%%%%%%%%%%%%%%%%%%%%%%%%%%%%%%%%%%%%%%%
\subsection{Angle Between Vectors}
\label{subsec:angle-between-vectors-reference}
%%%%%%%%%%%%%%%%%%%%%%%%%%%%%%%%%%%%%%%%%%%%%%%%%%%%%%%%%%%%

For nonzero vectors \(\mathbf{u},\mathbf{v}\),

\[
\boxed{
\cos\theta
=
\frac{
\mathbf{u}\cdot\mathbf{v}
}{
\|\mathbf{u}\|
\|\mathbf{v}\|
}.
}
\]

Therefore,

\[
\boxed{
\theta
=
\arccos
\left(
\frac{
\mathbf{u}\cdot\mathbf{v}
}{
\|\mathbf{u}\|
\|\mathbf{v}\|
}
\right).
}
\]

%%%%%%%%%%%%%%%%%%%%%%%%%%%%%%%%%%%%%%%%%%%%%%%%%%%%%%%%%%%%
\subsection{Orthogonal Vectors}
\label{subsec:orthogonal-vectors-reference}
%%%%%%%%%%%%%%%%%%%%%%%%%%%%%%%%%%%%%%%%%%%%%%%%%%%%%%%%%%%%

Two vectors are perpendicular when

\[
\boxed{
\mathbf{u}\cdot\mathbf{v}=0.
}
\]

%%%%%%%%%%%%%%%%%%%%%%%%%%%%%%%%%%%%%%%%%%%%%%%%%%%%%%%%%%%%
\subsection{Projection of One Vector Onto Another}
\label{subsec:vector-projection-reference}
%%%%%%%%%%%%%%%%%%%%%%%%%%%%%%%%%%%%%%%%%%%%%%%%%%%%%%%%%%%%

The projection of \(\mathbf{u}\) onto nonzero \(\mathbf{v}\) is

\[
\boxed{
\operatorname{proj}_{\mathbf{v}}\mathbf{u}
=
\frac{
\mathbf{u}\cdot\mathbf{v}
}{
\mathbf{v}\cdot\mathbf{v}
}
\mathbf{v}.
}
\]

%%%%%%%%%%%%%%%%%%%%%%%%%%%%%%%%%%%%%%%%%%%%%%%%%%%%%%%%%%%%
\subsection{Area of a Parallelogram Using Vectors}
\label{subsec:parallelogram-vector-area}
%%%%%%%%%%%%%%%%%%%%%%%%%%%%%%%%%%%%%%%%%%%%%%%%%%%%%%%%%%%%

For vectors \(\mathbf{u},\mathbf{v}\in\mathbb{R}^3\),

\[
\boxed{
A
=
\|\mathbf{u}\times\mathbf{v}\|.
}
\]

Since

\[
\|\mathbf{u}\times\mathbf{v}\|
=
\|\mathbf{u}\|
\|\mathbf{v}\|
\sin\theta,
\]

this agrees with the familiar formula \(ab\sin\theta\).

%%%%%%%%%%%%%%%%%%%%%%%%%%%%%%%%%%%%%%%%%%%%%%%%%%%%%%%%%%%%
\subsection{Triangle Area Using Vectors}
\label{subsec:triangle-vector-area}
%%%%%%%%%%%%%%%%%%%%%%%%%%%%%%%%%%%%%%%%%%%%%%%%%%%%%%%%%%%%

For vectors \(\mathbf{u},\mathbf{v}\) forming two sides of a triangle,

\[
\boxed{
A
=
\frac12
\|\mathbf{u}\times\mathbf{v}\|.
}
\]

%%%%%%%%%%%%%%%%%%%%%%%%%%%%%%%%%%%%%%%%%%%%%%%%%%%%%%%%%%%%
\subsection{Scalar Triple Product Volume}
\label{subsec:scalar-triple-product-volume}
%%%%%%%%%%%%%%%%%%%%%%%%%%%%%%%%%%%%%%%%%%%%%%%%%%%%%%%%%%%%

The volume of the parallelepiped generated by

\[
\mathbf{u},
\mathbf{v},
\mathbf{w}
\]

is

\[
\boxed{
V
=
|
\mathbf{u}\cdot
(\mathbf{v}\times\mathbf{w})
|.
}
\]

Equivalently,

\[
\boxed{
V
=
\left|
\det
\begin{bmatrix}
u_1 & u_2 & u_3\\
v_1 & v_2 & v_3\\
w_1 & w_2 & w_3
\end{bmatrix}
\right|.
}
\]

%%%%%%%%%%%%%%%%%%%%%%%%%%%%%%%%%%%%%%%%%%%%%%%%%%%%%%%%%%%%
\subsection{Equation of a Plane}
\label{subsec:plane-equation-reference}
%%%%%%%%%%%%%%%%%%%%%%%%%%%%%%%%%%%%%%%%%%%%%%%%%%%%%%%%%%%%

A plane with normal vector

\[
\mathbf{n}
=
(A,B,C)
\]

through point

\[
(x_0,y_0,z_0)
\]

has equation

\[
\boxed{
A(x-x_0)
+
B(y-y_0)
+
C(z-z_0)
=
0.
}
\]

Equivalent standard form is

\[
\boxed{
Ax+By+Cz=D.
}
\]

%%%%%%%%%%%%%%%%%%%%%%%%%%%%%%%%%%%%%%%%%%%%%%%%%%%%%%%%%%%%
\subsection{Distance From a Point to a Plane}
\label{subsec:point-plane-distance}
%%%%%%%%%%%%%%%%%%%%%%%%%%%%%%%%%%%%%%%%%%%%%%%%%%%%%%%%%%%%

For point

\[
(x_0,y_0,z_0)
\]

and plane

\[
Ax+By+Cz+D=0,
\]

distance is

\[
\boxed{
d
=
\frac{
|Ax_0+By_0+Cz_0+D|
}{
\sqrt{A^2+B^2+C^2}
}.
}
\]

%%%%%%%%%%%%%%%%%%%%%%%%%%%%%%%%%%%%%%%%%%%%%%%%%%%%%%%%%%%%
\subsection{Rectangular Prism}
\label{subsec:rectangular-prism-formulas}
%%%%%%%%%%%%%%%%%%%%%%%%%%%%%%%%%%%%%%%%%%%%%%%%%%%%%%%%%%%%

For side lengths \(l,w,h\),

\[
\boxed{
V=lwh.
}
\]

Surface area is

\[
\boxed{
S
=
2lw+2lh+2wh.
}
\]

Space diagonal is

\[
\boxed{
d
=
\sqrt{
l^2+w^2+h^2
}.
}
\]

%%%%%%%%%%%%%%%%%%%%%%%%%%%%%%%%%%%%%%%%%%%%%%%%%%%%%%%%%%%%
\subsection{Cube}
\label{subsec:cube-formulas}
%%%%%%%%%%%%%%%%%%%%%%%%%%%%%%%%%%%%%%%%%%%%%%%%%%%%%%%%%%%%

For side \(s\),

\[
\boxed{
V=s^3,
}
\]

\[
\boxed{
S=6s^2,
}
\]

and space diagonal

\[
\boxed{
d=s\sqrt3.
}
\]

%%%%%%%%%%%%%%%%%%%%%%%%%%%%%%%%%%%%%%%%%%%%%%%%%%%%%%%%%%%%
\subsection{General Prism}
\label{subsec:general-prism-formulas}
%%%%%%%%%%%%%%%%%%%%%%%%%%%%%%%%%%%%%%%%%%%%%%%%%%%%%%%%%%%%

For base area \(B\) and perpendicular height \(h\),

\[
\boxed{
V=Bh.
}
\]

For a right prism with base perimeter \(P\),

\[
\boxed{
S_{\mathrm{lateral}}
=
Ph.
}
\]

Total surface area is

\[
\boxed{
S_{\mathrm{total}}
=
Ph+2B.
}
\]

%%%%%%%%%%%%%%%%%%%%%%%%%%%%%%%%%%%%%%%%%%%%%%%%%%%%%%%%%%%%
\subsection{Cylinder}
\label{subsec:cylinder-formulas}
%%%%%%%%%%%%%%%%%%%%%%%%%%%%%%%%%%%%%%%%%%%%%%%%%%%%%%%%%%%%

For radius \(r\) and height \(h\),

\[
\boxed{
V
=
\pi r^2h.
}
\]

Lateral surface area is

\[
\boxed{
S_{\mathrm{lateral}}
=
2\pi rh.
}
\]

Total surface area is

\[
\boxed{
S_{\mathrm{total}}
=
2\pi rh+2\pi r^2.
}
\]

%%%%%%%%%%%%%%%%%%%%%%%%%%%%%%%%%%%%%%%%%%%%%%%%%%%%%%%%%%%%
\subsection{General Pyramid}
\label{subsec:pyramid-formulas}
%%%%%%%%%%%%%%%%%%%%%%%%%%%%%%%%%%%%%%%%%%%%%%%%%%%%%%%%%%%%

For base area \(B\) and perpendicular height \(h\),

\[
\boxed{
V
=
\frac13Bh.
}
\]

For a regular pyramid with base perimeter \(P\) and slant height \(\ell\),

\[
\boxed{
S_{\mathrm{lateral}}
=
\frac12P\ell.
}
\]

%%%%%%%%%%%%%%%%%%%%%%%%%%%%%%%%%%%%%%%%%%%%%%%%%%%%%%%%%%%%
\subsection{Cone}
\label{subsec:cone-formulas}
%%%%%%%%%%%%%%%%%%%%%%%%%%%%%%%%%%%%%%%%%%%%%%%%%%%%%%%%%%%%

For radius \(r\) and height \(h\),

\[
\boxed{
V
=
\frac13\pi r^2h.
}
\]

For a right circular cone, slant height is

\[
\boxed{
\ell
=
\sqrt{r^2+h^2}.
}
\]

Lateral surface area is

\[
\boxed{
S_{\mathrm{lateral}}
=
\pi r\ell.
}
\]

Total surface area is

\[
\boxed{
S_{\mathrm{total}}
=
\pi r\ell+\pi r^2.
}
\]

%%%%%%%%%%%%%%%%%%%%%%%%%%%%%%%%%%%%%%%%%%%%%%%%%%%%%%%%%%%%
\subsection{Sphere}
\label{subsec:sphere-formulas}
%%%%%%%%%%%%%%%%%%%%%%%%%%%%%%%%%%%%%%%%%%%%%%%%%%%%%%%%%%%%

For radius \(r\),

\[
\boxed{
V
=
\frac43\pi r^3.
}
\]

Surface area is

\[
\boxed{
S
=
4\pi r^2.
}
\]

%%%%%%%%%%%%%%%%%%%%%%%%%%%%%%%%%%%%%%%%%%%%%%%%%%%%%%%%%%%%
\subsection{Hemisphere}
\label{subsec:hemisphere-formulas}
%%%%%%%%%%%%%%%%%%%%%%%%%%%%%%%%%%%%%%%%%%%%%%%%%%%%%%%%%%%%

Volume is

\[
\boxed{
V
=
\frac23\pi r^3.
}
\]

Curved surface area is

\[
\boxed{
S_{\mathrm{curved}}
=
2\pi r^2.
}
\]

Including the circular base,

\[
\boxed{
S_{\mathrm{total}}
=
3\pi r^2.
}
\]

%%%%%%%%%%%%%%%%%%%%%%%%%%%%%%%%%%%%%%%%%%%%%%%%%%%%%%%%%%%%
\subsection{Frustum of a Cone}
\label{subsec:cone-frustum-formulas}
%%%%%%%%%%%%%%%%%%%%%%%%%%%%%%%%%%%%%%%%%%%%%%%%%%%%%%%%%%%%

For lower radius \(R\), upper radius \(r\), and height \(h\),

\[
\boxed{
V
=
\frac{\pi h}{3}
\left(
R^2+Rr+r^2
\right).
}
\]

The slant height is

\[
\boxed{
\ell
=
\sqrt{
h^2+(R-r)^2
}.
}
\]

Lateral area is

\[
\boxed{
S_{\mathrm{lateral}}
=
\pi(R+r)\ell.
}
\]

%%%%%%%%%%%%%%%%%%%%%%%%%%%%%%%%%%%%%%%%%%%%%%%%%%%%%%%%%%%%
\subsection{Frustum of a Pyramid}
\label{subsec:pyramid-frustum-volume}
%%%%%%%%%%%%%%%%%%%%%%%%%%%%%%%%%%%%%%%%%%%%%%%%%%%%%%%%%%%%

For base areas \(B_1,B_2\) and height \(h\),

\[
\boxed{
V
=
\frac{h}{3}
\left(
B_1
+
\sqrt{B_1B_2}
+
B_2
\right).
}
\]

%%%%%%%%%%%%%%%%%%%%%%%%%%%%%%%%%%%%%%%%%%%%%%%%%%%%%%%%%%%%
\subsection{Torus}
\label{subsec:torus-formulas}
%%%%%%%%%%%%%%%%%%%%%%%%%%%%%%%%%%%%%%%%%%%%%%%%%%%%%%%%%%%%

For major radius \(R\) and minor radius \(r\),

\[
\boxed{
V
=
2\pi^2Rr^2.
}
\]

Surface area is

\[
\boxed{
S
=
4\pi^2Rr.
}
\]

Here \(R\) is the distance from the center of the tube to the center of the
torus.

%%%%%%%%%%%%%%%%%%%%%%%%%%%%%%%%%%%%%%%%%%%%%%%%%%%%%%%%%%%%
\subsection{Regular Tetrahedron}
\label{subsec:regular-tetrahedron-formulas}
%%%%%%%%%%%%%%%%%%%%%%%%%%%%%%%%%%%%%%%%%%%%%%%%%%%%%%%%%%%%

For side length \(a\),

\[
\boxed{
S
=
\sqrt3\,a^2.
}
\]

Volume is

\[
\boxed{
V
=
\frac{a^3}{6\sqrt2}
=
\frac{\sqrt2}{12}a^3.
}
\]

Height is

\[
\boxed{
h
=
a\sqrt{\frac23}.
}
\]

%%%%%%%%%%%%%%%%%%%%%%%%%%%%%%%%%%%%%%%%%%%%%%%%%%%%%%%%%%%%
\subsection{Similarity}
\label{subsec:similarity-scaling-reference}
%%%%%%%%%%%%%%%%%%%%%%%%%%%%%%%%%%%%%%%%%%%%%%%%%%%%%%%%%%%%

If two figures are similar with length scale factor

\[
k,
\]

then corresponding lengths scale by

\[
\boxed{
k.
}
\]

Areas scale by

\[
\boxed{
k^2.
}
\]

Volumes scale by

\[
\boxed{
k^3.
}
\]

%%%%%%%%%%%%%%%%%%%%%%%%%%%%%%%%%%%%%%%%%%%%%%%%%%%%%%%%%%%%
\subsection{Worked Example: Scaling}
\label{subsec:worked-geometric-scaling}
%%%%%%%%%%%%%%%%%%%%%%%%%%%%%%%%%%%%%%%%%%%%%%%%%%%%%%%%%%%%

\begin{workedexamplebox}{Scaling a Sphere}

Suppose the radius of a sphere is doubled.

The radius scale factor is

\[
k=2.
\]

Surface area scales by

\[
k^2=4.
\]

Volume scales by

\[
k^3=8.
\]

\begin{answerbox}
\[
\boxed{
\text{surface area becomes }4\text{ times as large},
}
\]

\[
\boxed{
\text{volume becomes }8\text{ times as large}.
}
\]
\end{answerbox}

\end{workedexamplebox}

%%%%%%%%%%%%%%%%%%%%%%%%%%%%%%%%%%%%%%%%%%%%%%%%%%%%%%%%%%%%
\subsection{Law of Sines}
\label{subsec:law-of-sines-reference}
%%%%%%%%%%%%%%%%%%%%%%%%%%%%%%%%%%%%%%%%%%%%%%%%%%%%%%%%%%%%

For triangle sides \(a,b,c\) opposite angles \(A,B,C\),

\[
\boxed{
\frac{a}{\sin A}
=
\frac{b}{\sin B}
=
\frac{c}{\sin C}.
}
\]

If \(R\) is circumradius,

\[
\boxed{
\frac{a}{\sin A}
=
\frac{b}{\sin B}
=
\frac{c}{\sin C}
=
2R.
}
\]

%%%%%%%%%%%%%%%%%%%%%%%%%%%%%%%%%%%%%%%%%%%%%%%%%%%%%%%%%%%%
\subsection{Law of Cosines}
\label{subsec:law-of-cosines-reference}
%%%%%%%%%%%%%%%%%%%%%%%%%%%%%%%%%%%%%%%%%%%%%%%%%%%%%%%%%%%%

For a triangle,

\[
\boxed{
c^2
=
a^2+b^2-2ab\cos C.
}
\]

Similarly,

\[
\boxed{
a^2
=
b^2+c^2-2bc\cos A
}
\]

and

\[
\boxed{
b^2
=
a^2+c^2-2ac\cos B.
}
\]

The Pythagorean theorem is the special case \(C=90^\circ\).

%%%%%%%%%%%%%%%%%%%%%%%%%%%%%%%%%%%%%%%%%%%%%%%%%%%%%%%%%%%%
\subsection{Triangle Circumradius}
\label{subsec:triangle-circumradius}
%%%%%%%%%%%%%%%%%%%%%%%%%%%%%%%%%%%%%%%%%%%%%%%%%%%%%%%%%%%%

For side lengths \(a,b,c\) and area \(A\),

\[
\boxed{
R
=
\frac{abc}{4A}.
}
\]

Here \(R\) is the circumradius.

%%%%%%%%%%%%%%%%%%%%%%%%%%%%%%%%%%%%%%%%%%%%%%%%%%%%%%%%%%%%
\subsection{Triangle Inradius}
\label{subsec:triangle-inradius}
%%%%%%%%%%%%%%%%%%%%%%%%%%%%%%%%%%%%%%%%%%%%%%%%%%%%%%%%%%%%

If \(A\) is triangle area and

\[
s
=
\frac{a+b+c}{2}
\]

is semiperimeter, then inradius is

\[
\boxed{
r
=
\frac{A}{s}.
}
\]

Equivalently,

\[
\boxed{
A=rs.
}
\]

%%%%%%%%%%%%%%%%%%%%%%%%%%%%%%%%%%%%%%%%%%%%%%%%%%%%%%%%%%%%
\subsection{Centroid of a Triangle}
\label{subsec:triangle-centroid-reference}
%%%%%%%%%%%%%%%%%%%%%%%%%%%%%%%%%%%%%%%%%%%%%%%%%%%%%%%%%%%%

For vertices

\[
(x_1,y_1),
\quad
(x_2,y_2),
\quad
(x_3,y_3),
\]

the centroid is

\[
\boxed{
G
=
\left(
\frac{x_1+x_2+x_3}{3},
\frac{y_1+y_2+y_3}{3}
\right).
}
\]

%%%%%%%%%%%%%%%%%%%%%%%%%%%%%%%%%%%%%%%%%%%%%%%%%%%%%%%%%%%%
\subsection{Centroid of a Tetrahedron}
\label{subsec:tetrahedron-centroid-reference}
%%%%%%%%%%%%%%%%%%%%%%%%%%%%%%%%%%%%%%%%%%%%%%%%%%%%%%%%%%%%

For vertices

\[
\mathbf{v}_1,
\mathbf{v}_2,
\mathbf{v}_3,
\mathbf{v}_4,
\]

the centroid is

\[
\boxed{
\mathbf{G}
=
\frac14
\left(
\mathbf{v}_1+
\mathbf{v}_2+
\mathbf{v}_3+
\mathbf{v}_4
\right).
}
\]

%%%%%%%%%%%%%%%%%%%%%%%%%%%%%%%%%%%%%%%%%%%%%%%%%%%%%%%%%%%%
\subsection{Polar Coordinates}
\label{subsec:polar-coordinate-formulas}
%%%%%%%%%%%%%%%%%%%%%%%%%%%%%%%%%%%%%%%%%%%%%%%%%%%%%%%%%%%%

The coordinate transformations are

\[
\boxed{
x=r\cos\theta,
}
\]

\[
\boxed{
y=r\sin\theta.
}
\]

Conversely,

\[
\boxed{
r
=
\sqrt{x^2+y^2}.
}
\]

For \(x\neq0\),

\[
\tan\theta
=
\frac{y}{x},
\]

with quadrant information used to determine the correct angle.

%%%%%%%%%%%%%%%%%%%%%%%%%%%%%%%%%%%%%%%%%%%%%%%%%%%%%%%%%%%%
\subsection{Cylindrical Coordinates}
\label{subsec:cylindrical-coordinate-formulas}
%%%%%%%%%%%%%%%%%%%%%%%%%%%%%%%%%%%%%%%%%%%%%%%%%%%%%%%%%%%%

Cylindrical coordinates satisfy

\[
\boxed{
x=r\cos\theta,
}
\]

\[
\boxed{
y=r\sin\theta,
}
\]

\[
\boxed{
z=z.
}
\]

Thus

\[
\boxed{
r^2=x^2+y^2.
}
\]

%%%%%%%%%%%%%%%%%%%%%%%%%%%%%%%%%%%%%%%%%%%%%%%%%%%%%%%%%%%%
\subsection{Spherical Coordinates}
\label{subsec:spherical-coordinate-formulas}
%%%%%%%%%%%%%%%%%%%%%%%%%%%%%%%%%%%%%%%%%%%%%%%%%%%%%%%%%%%%

Using the convention where \(\rho\) is radial distance,
\(\theta\) is azimuth, and \(\phi\) is polar angle,

\[
\boxed{
x
=
\rho\sin\phi\cos\theta,
}
\]

\[
\boxed{
y
=
\rho\sin\phi\sin\theta,
}
\]

\[
\boxed{
z
=
\rho\cos\phi.
}
\]

Also,

\[
\boxed{
\rho
=
\sqrt{x^2+y^2+z^2}.
}
\]

Spherical-coordinate conventions vary, so the angle definitions should
always be stated.

%%%%%%%%%%%%%%%%%%%%%%%%%%%%%%%%%%%%%%%%%%%%%%%%%%%%%%%%%%%%
\subsection{Parametric Line}
\label{subsec:parametric-line-reference}
%%%%%%%%%%%%%%%%%%%%%%%%%%%%%%%%%%%%%%%%%%%%%%%%%%%%%%%%%%%%

A line through point \(\mathbf{r}_0\) with direction \(\mathbf{v}\) is

\[
\boxed{
\mathbf{r}(t)
=
\mathbf{r}_0+t\mathbf{v}.
}
\]

In coordinates,

\[
\boxed{
x=x_0+at,
\qquad
y=y_0+bt,
\qquad
z=z_0+ct.
}
\]

%%%%%%%%%%%%%%%%%%%%%%%%%%%%%%%%%%%%%%%%%%%%%%%%%%%%%%%%%%%%
\subsection{Circle Parametrization}
\label{subsec:circle-parametrization-reference}
%%%%%%%%%%%%%%%%%%%%%%%%%%%%%%%%%%%%%%%%%%%%%%%%%%%%%%%%%%%%

A circle of radius \(r\) centered at \((h,k)\) can be parameterized by

\[
\boxed{
x(t)
=
h+r\cos t,
}
\]

\[
\boxed{
y(t)
=
k+r\sin t.
}
\]

For one full traversal,

\[
\boxed{
0\leq t\leq2\pi.
}
\]

%%%%%%%%%%%%%%%%%%%%%%%%%%%%%%%%%%%%%%%%%%%%%%%%%%%%%%%%%%%%
\subsection{Ellipse Parametrization}
\label{subsec:ellipse-parametrization-reference}
%%%%%%%%%%%%%%%%%%%%%%%%%%%%%%%%%%%%%%%%%%%%%%%%%%%%%%%%%%%%

An ellipse centered at \((h,k)\) with semiaxes \(a,b\) has parametrization

\[
\boxed{
x(t)
=
h+a\cos t,
}
\]

\[
\boxed{
y(t)
=
k+b\sin t.
}
\]

%%%%%%%%%%%%%%%%%%%%%%%%%%%%%%%%%%%%%%%%%%%%%%%%%%%%%%%%%%%%
\subsection{Arc Length of a Plane Parametric Curve}
\label{subsec:parametric-arc-length-reference}
%%%%%%%%%%%%%%%%%%%%%%%%%%%%%%%%%%%%%%%%%%%%%%%%%%%%%%%%%%%%

For

\[
x=x(t),
\qquad
y=y(t),
\]

with \(a\leq t\leq b\),

\[
\boxed{
L
=
\int_a^b
\sqrt{
\left(
\frac{dx}{dt}
\right)^2
+
\left(
\frac{dy}{dt}
\right)^2
}
\,dt.
}
\]

%%%%%%%%%%%%%%%%%%%%%%%%%%%%%%%%%%%%%%%%%%%%%%%%%%%%%%%%%%%%
\subsection{Arc Length of a Spatial Curve}
\label{subsec:spatial-arc-length-reference}
%%%%%%%%%%%%%%%%%%%%%%%%%%%%%%%%%%%%%%%%%%%%%%%%%%%%%%%%%%%%

For

\[
\mathbf{r}(t)
=
(x(t),y(t),z(t)),
\]

\[
\boxed{
L
=
\int_a^b
\|\mathbf{r}'(t)\|
\,dt.
}
\]

Thus

\[
\boxed{
L
=
\int_a^b
\sqrt{
(x')^2+(y')^2+(z')^2
}
\,dt.
}
\]

%%%%%%%%%%%%%%%%%%%%%%%%%%%%%%%%%%%%%%%%%%%%%%%%%%%%%%%%%%%%
\subsection{Surface Area of Revolution About the \(x\)-Axis}
\label{subsec:surface-area-revolution-x}
%%%%%%%%%%%%%%%%%%%%%%%%%%%%%%%%%%%%%%%%%%%%%%%%%%%%%%%%%%%%

For \(y=f(x)\geq0\) on \([a,b]\),

\[
\boxed{
S
=
2\pi
\int_a^b
f(x)
\sqrt{
1+[f'(x)]^2
}
\,dx.
}
\]

%%%%%%%%%%%%%%%%%%%%%%%%%%%%%%%%%%%%%%%%%%%%%%%%%%%%%%%%%%%%
\subsection{Surface Area of Revolution About the \(y\)-Axis}
\label{subsec:surface-area-revolution-y}
%%%%%%%%%%%%%%%%%%%%%%%%%%%%%%%%%%%%%%%%%%%%%%%%%%%%%%%%%%%%

For \(x\geq0\),

\[
\boxed{
S
=
2\pi
\int_a^b
x
\sqrt{
1+[f'(x)]^2
}
\,dx.
}
\]

Here \(y=f(x)\).

%%%%%%%%%%%%%%%%%%%%%%%%%%%%%%%%%%%%%%%%%%%%%%%%%%%%%%%%%%%%
\subsection{Volume by Disks and Washers}
\label{subsec:disk-washer-volume-reference}
%%%%%%%%%%%%%%%%%%%%%%%%%%%%%%%%%%%%%%%%%%%%%%%%%%%%%%%%%%%%

For rotation about the \(x\)-axis,

\[
\boxed{
V
=
\pi
\int_a^b
\left[
R(x)^2-r(x)^2
\right]
dx.
}
\]

Here \(R\) is outer radius and \(r\) is inner radius.

%%%%%%%%%%%%%%%%%%%%%%%%%%%%%%%%%%%%%%%%%%%%%%%%%%%%%%%%%%%%
\subsection{Volume by Cylindrical Shells}
\label{subsec:shell-volume-reference}
%%%%%%%%%%%%%%%%%%%%%%%%%%%%%%%%%%%%%%%%%%%%%%%%%%%%%%%%%%%%

For shells of radius \(r(x)\) and height \(h(x)\),

\[
\boxed{
V
=
2\pi
\int_a^b
r(x)h(x)\,dx.
}
\]

%%%%%%%%%%%%%%%%%%%%%%%%%%%%%%%%%%%%%%%%%%%%%%%%%%%%%%%%%%%%
\subsection{Area in Polar Coordinates}
\label{subsec:polar-area-reference}
%%%%%%%%%%%%%%%%%%%%%%%%%%%%%%%%%%%%%%%%%%%%%%%%%%%%%%%%%%%%

For polar curve

\[
r=f(\theta),
\]

area between \(\theta=\alpha\) and \(\theta=\beta\) is

\[
\boxed{
A
=
\frac12
\int_\alpha^\beta
r^2\,d\theta.
}
\]

%%%%%%%%%%%%%%%%%%%%%%%%%%%%%%%%%%%%%%%%%%%%%%%%%%%%%%%%%%%%
\subsection{Polar Arc Length}
\label{subsec:polar-arc-length-reference}
%%%%%%%%%%%%%%%%%%%%%%%%%%%%%%%%%%%%%%%%%%%%%%%%%%%%%%%%%%%%

For

\[
r=r(\theta),
\]

\[
\boxed{
L
=
\int_\alpha^\beta
\sqrt{
r^2
+
\left(
\frac{dr}{d\theta}
\right)^2
}
\,d\theta.
}
\]

%%%%%%%%%%%%%%%%%%%%%%%%%%%%%%%%%%%%%%%%%%%%%%%%%%%%%%%%%%%%
\subsection{Curvature of a Plane Curve}
\label{subsec:plane-curvature-reference}
%%%%%%%%%%%%%%%%%%%%%%%%%%%%%%%%%%%%%%%%%%%%%%%%%%%%%%%%%%%%

For \(y=f(x)\),

\[
\boxed{
\kappa
=
\frac{
|f''(x)|
}{
\left(
1+[f'(x)]^2
\right)^{3/2}
}.
}
\]

The radius of curvature is

\[
\boxed{
R
=
\frac1\kappa
}
\]

when \(\kappa\neq0\).

%%%%%%%%%%%%%%%%%%%%%%%%%%%%%%%%%%%%%%%%%%%%%%%%%%%%%%%%%%%%
\subsection{Curvature of a Parametric Plane Curve}
\label{subsec:parametric-curvature-reference}
%%%%%%%%%%%%%%%%%%%%%%%%%%%%%%%%%%%%%%%%%%%%%%%%%%%%%%%%%%%%

For

\[
\mathbf{r}(t)
=
(x(t),y(t)),
\]

\[
\boxed{
\kappa
=
\frac{
|x'y''-y'x''|
}{
\left[
(x')^2+(y')^2
\right]^{3/2}
}.
}
\]

%%%%%%%%%%%%%%%%%%%%%%%%%%%%%%%%%%%%%%%%%%%%%%%%%%%%%%%%%%%%
\subsection{Curvature of a Space Curve}
\label{subsec:space-curve-curvature-reference}
%%%%%%%%%%%%%%%%%%%%%%%%%%%%%%%%%%%%%%%%%%%%%%%%%%%%%%%%%%%%

For a sufficiently regular curve,

\[
\boxed{
\kappa
=
\frac{
\|\mathbf{r}'\times\mathbf{r}''\|
}{
\|\mathbf{r}'\|^3
}.
}
\]

%%%%%%%%%%%%%%%%%%%%%%%%%%%%%%%%%%%%%%%%%%%%%%%%%%%%%%%%%%%%
\subsection{Quick Plane Geometry Table}
\label{subsec:quick-plane-geometry-table}
%%%%%%%%%%%%%%%%%%%%%%%%%%%%%%%%%%%%%%%%%%%%%%%%%%%%%%%%%%%%

\begin{center}
\begin{tabular}{c|c|c}
\textbf{Shape}
&
\textbf{Area}
&
\textbf{Perimeter / Circumference}
\\ \hline

Square
&
\(s^2\)
&
\(4s\)
\\

Rectangle
&
\(lw\)
&
\(2(l+w)\)
\\

Triangle
&
\(\frac12bh\)
&
\(a+b+c\)
\\

Parallelogram
&
\(bh\)
&
\(2(a+b)\)
\\

Trapezoid
&
\(\frac12(b_1+b_2)h\)
&
sum of sides
\\

Rhombus
&
\(\frac12d_1d_2\)
&
\(4s\)
\\

Circle
&
\(\pi r^2\)
&
\(2\pi r\)
\end{tabular}
\end{center}

%%%%%%%%%%%%%%%%%%%%%%%%%%%%%%%%%%%%%%%%%%%%%%%%%%%%%%%%%%%%
\subsection{Quick Solid Geometry Table}
\label{subsec:quick-solid-geometry-table}
%%%%%%%%%%%%%%%%%%%%%%%%%%%%%%%%%%%%%%%%%%%%%%%%%%%%%%%%%%%%

\begin{center}
\begin{tabular}{c|c|c}
\textbf{Solid}
&
\textbf{Volume}
&
\textbf{Surface Area}
\\ \hline

Cube
&
\(s^3\)
&
\(6s^2\)
\\

Rectangular prism
&
\(lwh\)
&
\(2lw+2lh+2wh\)
\\

Cylinder
&
\(\pi r^2h\)
&
\(2\pi rh+2\pi r^2\)
\\

Cone
&
\(\frac13\pi r^2h\)
&
\(\pi r\ell+\pi r^2\)
\\

Sphere
&
\(\frac43\pi r^3\)
&
\(4\pi r^2\)
\\

Pyramid
&
\(\frac13Bh\)
&
depends on base and slant height
\end{tabular}
\end{center}

%%%%%%%%%%%%%%%%%%%%%%%%%%%%%%%%%%%%%%%%%%%%%%%%%%%%%%%%%%%%
\subsection{Quick Coordinate Geometry Table}
\label{subsec:quick-coordinate-geometry-table}
%%%%%%%%%%%%%%%%%%%%%%%%%%%%%%%%%%%%%%%%%%%%%%%%%%%%%%%%%%%%

\[
\boxed{
d
=
\sqrt{
(x_2-x_1)^2+(y_2-y_1)^2
}
}
\]

\[
\boxed{
M
=
\left(
\frac{x_1+x_2}{2},
\frac{y_1+y_2}{2}
\right)
}
\]

\[
\boxed{
m
=
\frac{
y_2-y_1
}{
x_2-x_1
}
}
\]

\[
\boxed{
y-y_1
=
m(x-x_1)
}
\]

\[
\boxed{
(x-h)^2+(y-k)^2=r^2
}
\]

%%%%%%%%%%%%%%%%%%%%%%%%%%%%%%%%%%%%%%%%%%%%%%%%%%%%%%%%%%%%
\subsection{Common Errors}
\label{subsec:geometric-formulas-common-errors}
%%%%%%%%%%%%%%%%%%%%%%%%%%%%%%%%%%%%%%%%%%%%%%%%%%%%%%%%%%%%

\subsubsection{Using Diameter Where Radius Is Required}

Circle area is

\[
\boxed{
A=\pi r^2,
}
\]

not

\[
\pi d^2.
\]

Since

\[
r=\frac d2,
\]

one may write

\[
\boxed{
A
=
\frac{\pi d^2}{4}.
}
\]

%%%%%%%%%%%%%%%%%%%%%%%%%%%%%%%%%%%%%%%%%%%%%%%%%%%%%%%%%%%%
\subsubsection{Forgetting to Square the Radius}

Circumference is linear in radius:

\[
C=2\pi r.
\]

Area is quadratic:

\[
A=\pi r^2.
\]

%%%%%%%%%%%%%%%%%%%%%%%%%%%%%%%%%%%%%%%%%%%%%%%%%%%%%%%%%%%%
\subsubsection{Confusing Area and Perimeter}

Perimeter has units of length.

Area has units of length squared.

For example,

\[
\boxed{
P
\text{ has units }L,
}
\]

while

\[
\boxed{
A
\text{ has units }L^2.
}
\]

%%%%%%%%%%%%%%%%%%%%%%%%%%%%%%%%%%%%%%%%%%%%%%%%%%%%%%%%%%%%
\subsubsection{Confusing Surface Area and Volume}

Surface area scales as

\[
\boxed{
L^2,
}
\]

while volume scales as

\[
\boxed{
L^3.
}
\]

Their units should never be interchanged.

%%%%%%%%%%%%%%%%%%%%%%%%%%%%%%%%%%%%%%%%%%%%%%%%%%%%%%%%%%%%
\subsubsection{Using Slant Height as Perpendicular Height}

Cone volume uses perpendicular height:

\[
\boxed{
V
=
\frac13\pi r^2h.
}
\]

Surface area uses slant height:

\[
\boxed{
S_{\mathrm{lateral}}
=
\pi r\ell.
}
\]

%%%%%%%%%%%%%%%%%%%%%%%%%%%%%%%%%%%%%%%%%%%%%%%%%%%%%%%%%%%%
\subsubsection{Using Degrees in \(s=r\theta\)}

The formula

\[
\boxed{
s=r\theta
}
\]

requires \(\theta\) in radians.

If degrees are used, convert first or use the degree-based fraction of a full
circle.

%%%%%%%%%%%%%%%%%%%%%%%%%%%%%%%%%%%%%%%%%%%%%%%%%%%%%%%%%%%%
\subsubsection{Using Degrees in \(\frac12r^2\theta\)}

Similarly,

\[
\boxed{
A_{\mathrm{sector}}
=
\frac12r^2\theta
}
\]

requires radians.

%%%%%%%%%%%%%%%%%%%%%%%%%%%%%%%%%%%%%%%%%%%%%%%%%%%%%%%%%%%%
\subsubsection{Forgetting Absolute Value in Coordinate Area}

The determinant or shoelace expression may be negative depending on vertex
orientation.

Geometric area is nonnegative:

\[
\boxed{
A
=
\frac12|\text{signed determinant}|.
}
\]

%%%%%%%%%%%%%%%%%%%%%%%%%%%%%%%%%%%%%%%%%%%%%%%%%%%%%%%%%%%%
\subsubsection{Using the Pythagorean Theorem on a Nonright Triangle}

The formula

\[
a^2+b^2=c^2
\]

requires a right angle.

For a general triangle use the Law of Cosines.

%%%%%%%%%%%%%%%%%%%%%%%%%%%%%%%%%%%%%%%%%%%%%%%%%%%%%%%%%%%%
\subsubsection{Mixing Semiperimeter and Perimeter in Heron's Formula}

Heron's formula uses

\[
\boxed{
s
=
\frac{a+b+c}{2},
}
\]

not the full perimeter.

%%%%%%%%%%%%%%%%%%%%%%%%%%%%%%%%%%%%%%%%%%%%%%%%%%%%%%%%%%%%
\subsubsection{Using the Wrong Hyperbola Relation}

For an ellipse,

\[
\boxed{
c^2=a^2-b^2.
}
\]

For a hyperbola,

\[
\boxed{
c^2=a^2+b^2.
}
\]

%%%%%%%%%%%%%%%%%%%%%%%%%%%%%%%%%%%%%%%%%%%%%%%%%%%%%%%%%%%%
\subsubsection{Assuming All Spherical Coordinate Conventions Match}

Some texts interchange the meanings of \(\theta\) and \(\phi\).

Always check the stated convention.

%%%%%%%%%%%%%%%%%%%%%%%%%%%%%%%%%%%%%%%%%%%%%%%%%%%%%%%%%%%%
\subsubsection{Ignoring Scale Dimensions}

Doubling a length does not merely double area or volume.

If \(k=2\),

\[
\boxed{
A
\mapsto
4A,
}
\]

and

\[
\boxed{
V
\mapsto
8V.
}
\]

%%%%%%%%%%%%%%%%%%%%%%%%%%%%%%%%%%%%%%%%%%%%%%%%%%%%%%%%%%%%
\subsection{Exercises}
\label{subsec:geometric-formulas-exercises}
%%%%%%%%%%%%%%%%%%%%%%%%%%%%%%%%%%%%%%%%%%%%%%%%%%%%%%%%%%%%

\begin{exercise}[1]
Convert \(150^\circ\) to radians.
\end{exercise}

\begin{exercise}[1]
Convert \(7\pi/6\) radians to degrees.
\end{exercise}

\begin{exercise}[1]
Find the hypotenuse of a right triangle with legs \(8\) and \(15\).
\end{exercise}

\begin{exercise}[1]
Find the area of a triangle with base \(12\) and height \(7\).
\end{exercise}

\begin{exercise}[1]
Find the perimeter and area of a rectangle with dimensions \(5\) by \(9\).
\end{exercise}

\begin{exercise}[1]
Find the diagonal of a square with side \(10\).
\end{exercise}

\begin{exercise}[1]
Find the circumference and area of a circle with radius \(6\).
\end{exercise}

\begin{exercise}[2]
Find the arc length subtended by angle \(\pi/3\) in a circle of radius \(9\).
\end{exercise}

\begin{exercise}[2]
Find the area of a sector with radius \(4\) and central angle \(5\pi/6\).
\end{exercise}

\begin{exercise}[2]
Use Heron's formula for a triangle with sides \(5,5,6\).
\end{exercise}

\begin{exercise}[2]
Use the Law of Cosines to find the third side of a triangle with sides
\(7\), \(10\), and included angle \(60^\circ\).
\end{exercise}

\begin{exercise}[2]
Use the Law of Sines to find a missing side in a nonright triangle.
\end{exercise}

\begin{exercise}[2]
Find the area of a regular hexagon with side length \(4\).
\end{exercise}

\begin{exercise}[2]
Find the distance between

\[
(2,-1)
\]

and

\[
(8,7).
\]
\end{exercise}

\begin{exercise}[2]
Find the midpoint between

\[
(-4,3)
\]

and

\[
(6,-5).
\]
\end{exercise}

\begin{exercise}[2]
Find the slope of the line through

\[
(1,2)
\]

and

\[
(5,10).
\]
\end{exercise}

\begin{exercise}[2]
Find the equation of a line with slope \(3\) through \((2,-1)\).
\end{exercise}

\begin{exercise}[2]
Find the distance from \((3,4)\) to the line

\[
2x-y+5=0.
\]
\end{exercise}

\begin{exercise}[2]
Write the equation of a circle centered at \((3,-2)\) with radius \(5\).
\end{exercise}

\begin{exercise}[2]
Find the area of the triangle with vertices

\[
(0,0),
\quad
(4,0),
\quad
(1,3).
\]
\end{exercise}

\begin{exercise}[3]
Use the shoelace formula to compute the area of a quadrilateral.
\end{exercise}

\begin{exercise}[2]
Find the volume and surface area of a cube with side \(6\).
\end{exercise}

\begin{exercise}[2]
Find the volume of a cylinder with radius \(3\) and height \(10\).
\end{exercise}

\begin{exercise}[2]
Find the surface area of the same cylinder.
\end{exercise}

\begin{exercise}[2]
Find the volume of a cone with radius \(5\) and height \(12\).
\end{exercise}

\begin{exercise}[2]
Find the slant height and surface area of the same cone.
\end{exercise}

\begin{exercise}[2]
Find the volume and surface area of a sphere with radius \(7\).
\end{exercise}

\begin{exercise}[2]
A solid is enlarged by scale factor \(3\). Determine the area and volume
scale factors.
\end{exercise}

\begin{exercise}[2]
Compute the angle between

\[
\mathbf{u}=(1,2)
\]

and

\[
\mathbf{v}=(3,-1).
\]
\end{exercise}

\begin{exercise}[2]
Determine whether

\[
(2,3,-1)
\]

and

\[
(1,-1,-1)
\]

are orthogonal.
\end{exercise}

\begin{exercise}[2]
Compute the projection of

\[
(3,4)
\]

onto

\[
(1,0).
\]
\end{exercise}

\begin{exercise}[3]
Find the area of the parallelogram generated by

\[
\mathbf{u}=(1,2,3),
\qquad
\mathbf{v}=(2,-1,1).
\]
\end{exercise}

\begin{exercise}[3]
Find the volume of the parallelepiped generated by three specified vectors.
\end{exercise}

\begin{exercise}[2]
Find the equation of the plane through \((1,2,3)\) with normal
\((2,-1,4)\).
\end{exercise}

\begin{exercise}[2]
Find the distance from a point to a specified plane.
\end{exercise}

\begin{exercise}[2]
Convert the point

\[
(3,3)
\]

from rectangular coordinates to polar coordinates.
\end{exercise}

\begin{exercise}[2]
Convert

\[
(r,\theta)
=
\left(
4,\frac{\pi}{6}
\right)
\]

to rectangular coordinates.
\end{exercise}

\begin{exercise}[2]
Parameterize a circle of radius \(5\) centered at \((2,1)\).
\end{exercise}

\begin{exercise}[2]
Parameterize an ellipse with semiaxes \(4\) and \(2\).
\end{exercise}

\begin{exercise}[3]
Compute arc length of

\[
\mathbf{r}(t)
=
(t,t^2),
\qquad
0\leq t\leq1.
\]
\end{exercise}

\begin{exercise}[3]
Find polar area enclosed by

\[
r=2
\]

over one full revolution.
\end{exercise}

\begin{exercise}[3]
Verify that the polar-area formula reproduces \(\pi r^2\) for a circle.
\end{exercise}

\begin{exercise}[3]
Use the disk method to derive the volume of a sphere.
\end{exercise}

\begin{exercise}[3]
Use cylindrical shells to compute the volume of a specified solid of
revolution.
\end{exercise}

\begin{exercise}[3]
Find the curvature of

\[
y=x^2
\]

at \(x=0\).
\end{exercise}

\begin{exercise}[3]
Find the centroid of a triangle from its coordinates.
\end{exercise}

\begin{exercise}[3]
Derive the circumradius relation

\[
R=\frac{abc}{4A}.
\]
\end{exercise}

\begin{exercise}[3]
Show that the area of a triangle is

\[
A=rs,
\]

where \(r\) is inradius and \(s\) is semiperimeter.
\end{exercise}

\begin{exercise}[4]
Derive Heron's formula.
\end{exercise}

\begin{exercise}[4]
Derive the regular polygon area formula

\[
A
=
\frac{ns^2}{4\tan(\pi/n)}.
\]
\end{exercise}

\begin{exercise}[4]
Derive the cone volume formula using integration.
\end{exercise}

\begin{exercise}[4]
Derive the sphere surface-area formula.
\end{exercise}

\begin{exercise}[4]
Derive the sphere volume formula.
\end{exercise}

\begin{exercise}[4]
Derive the conical frustum volume formula.
\end{exercise}

\begin{exercise}[4]
Derive the shoelace formula using determinants.
\end{exercise}

\begin{exercise}[4]
Prove the scalar triple product gives parallelepiped volume.
\end{exercise}

\begin{exercise}[4]
Derive the formula for distance from a point to a plane.
\end{exercise}

\begin{exercise}[4]
Derive the curvature formula for a plane graph \(y=f(x)\).
\end{exercise}

%%%%%%%%%%%%%%%%%%%%%%%%%%%%%%%%%%%%%%%%%%%%%%%%%%%%%%%%%%%%
\subsection{Conceptual Summary}
\label{subsec:geometric-formulas-summary}
%%%%%%%%%%%%%%%%%%%%%%%%%%%%%%%%%%%%%%%%%%%%%%%%%%%%%%%%%%%%

Geometry connects algebraic measurements with spatial structure.

For right triangles,

\[
\boxed{
a^2+b^2=c^2.
}
\]

For a triangle,

\[
\boxed{
A
=
\frac12bh.
}
\]

For circles,

\[
\boxed{
C=2\pi r
}
\]

and

\[
\boxed{
A=\pi r^2.
}
\]

For spheres,

\[
\boxed{
S=4\pi r^2
}
\]

and

\[
\boxed{
V=\frac43\pi r^3.
}
\]

Coordinate geometry uses

\[
\boxed{
d
=
\sqrt{
(x_2-x_1)^2+(y_2-y_1)^2
}
}
\]

and

\[
\boxed{
m
=
\frac{y_2-y_1}{x_2-x_1}.
}
\]

Vector geometry extends these formulas through dot products, cross products,
and determinants.

Scaling by a length factor \(k\) changes

\[
\boxed{
\text{length by }k,
}
\]

\[
\boxed{
\text{area by }k^2,
}
\]

and

\[
\boxed{
\text{volume by }k^3.
}
\]

The most important practical principles are:

\[
\boxed{
\text{use consistent units},
}
\]

\[
\boxed{
\text{distinguish length, area, and volume},
}
\]

and

\[
\boxed{
\text{identify which geometric measurements each formula requires}.
}
\]

%%%%%%%%%%%%%%%%%%%%%%%%%%%%%%%%%%%%%%%%%%%%%%%%%%%%%%%%%%%%
\subsection{Section Connections}
\label{subsec:geometric-formulas-connections}
%%%%%%%%%%%%%%%%%%%%%%%%%%%%%%%%%%%%%%%%%%%%%%%%%%%%%%%%%%%%

\chapterconnection{
Geometric Formulas consolidates the Euclidean, coordinate, vector, analytic,
and solid-geometry relationships used throughout the textbook. Distance,
slope, midpoint, circle, and conic formulas connect with Algebra and
Precalculus; triangle formulas connect with Trigonometry; vector area and
volume formulas connect with Linear Algebra and Vector Calculus; arc length,
surface area, and volumes of revolution connect with Calculus; curvature
connects with Differential Geometry; and coordinate transformations connect
with multivariable analysis and mathematical physics. The next reference
section, Trigonometric Identities, collects reciprocal, quotient,
Pythagorean, cofunction, sum-and-difference, double-angle, half-angle,
product-to-sum, sum-to-product, and inverse-trigonometric relationships.
}

%%%%%%%%%%%%%%%%%%%%%%%%%%%%%%%%%%%%%%%%%%%%%%%%%%%%%%%%%%%%
% SECTION SOURCE TRACKING
%%%%%%%%%%%%%%%%%%%%%%%%%%%%%%%%%%%%%%%%%%%%%%%%%%%%%%%%%%%%

%------------------------------------------------------------
% 14.6 Geometric Formulas
%------------------------------------------------------------
%
% Source basis:
%   - Standard Euclidean geometry.
%   - Standard coordinate and analytic geometry.
%   - Standard vector geometry and elementary calculus geometry.
%
% Used for:
%   - angular units;
%   - degree-radian conversion;
%   - triangles;
%   - polygon angle sums;
%   - Pythagorean theorem;
%   - special right triangles;
%   - triangle area;
%   - Heron's formula;
%   - rectangles and squares;
%   - parallelograms, rhombi, trapezoids, and kites;
%   - regular polygons;
%   - circles;
%   - arcs, sectors, chords, and segments;
%   - annuli;
%   - conic sections;
%   - coordinate distances;
%   - midpoints;
%   - lines and slopes;
%   - point-to-line distance;
%   - circles and spheres in coordinates;
%   - shoelace formula;
%   - vector geometry;
%   - cross products and triple products;
%   - planes;
%   - prisms;
%   - cylinders;
%   - pyramids;
%   - cones;
%   - spheres and hemispheres;
%   - frustums;
%   - tori;
%   - tetrahedra;
%   - similarity and scaling;
%   - Law of Sines;
%   - Law of Cosines;
%   - circumradius and inradius;
%   - centroids;
%   - polar, cylindrical, and spherical coordinates;
%   - parametrizations;
%   - arc length;
%   - surface area of revolution;
%   - volumes of revolution;
%   - polar area;
%   - curvature;
%   - common errors;
%   - applications and exercises.
%
% Notes:
%   - Trigonometric Identities is developed next in Section 14.7.
%   - Limit Laws follows in Section 14.8.
%
%%%%%%%%%%%%%%%%%%%%%%%%%%%%%%%%%%%%%%%%%%%%%%%%%%%%%%%%%%%%